% !TEX root = ../main.tex
%% Sujet type bac n° 5 — sujet inédit OptiTAB (valeurs vérifiées : verifications/sujets_bac.py)

\sujetBac{5}{{La pile d'une aide auditive : piles, sens d'évolution, capacité}/8,{Le saut en parachute : deuxième loi de Newton, méthode d'Euler}/7,{Le diamètre d'un cheveu : diffraction, incertitudes}/5}

\begin{exoBac}{8}{La pile d'une aide auditive}
Les aides auditives fonctionnent avec de petites piles bouton « zinc-air ». Avant la première utilisation, on retire une
pastille adhésive qui obture de minuscules trous : l'air entre alors dans la pile, qui commence à fonctionner.

\partieExo{1}{Le fonctionnement d'une pile}
\begin{docu}{1}{Une pile Daniell au laboratoire}
Une lame de zinc plonge dans une solution de sulfate de zinc, $\ce{Zn^{2+}}+\ce{SO4^{2-}}$, à \qty{0,10}{mol.L^{-1}} ; une lame de cuivre
plonge dans une solution de sulfate de cuivre, $\ce{Cu^{2+}}+\ce{SO4^{2-}}$, à \qty{0,10}{mol.L^{-1}}. Un pont salin relie les deux solutions.
La réaction $\ce{Cu^{2+} + Zn <=> Cu + Zn^{2+}}$ a pour constante d'équilibre $K=\num{1.9e37}$ à \qty{25}{\degreeCelsius}.
\end{docu}
\begin{questions}
\item Calculer le quotient de réaction initial et en déduire le sens d'évolution spontanée du système.
\item Indiquer la polarité de la pile et le sens de circulation des électrons dans le circuit extérieur. Quel est le rôle
du pont salin ?
\item Pourquoi la tension de la pile diminue-t-elle au cours de son utilisation ? Que vaut-elle lorsque la pile est usée ?
\end{questions}

\partieExo{2}{La pile zinc-air}
\begin{docu}{2}{Caractéristiques de la pile}
Demi-équations, en milieu basique : $\ce{Zn + 2HO- -> ZnO + H2O + 2e-}$ et $\ce{O2 + 2H2O + 4e- -> 4HO-}$. Tension : \qty{1,4}{V} ;
capacité : \qty{620}{mA.h} ; masse : \qty{1,9}{g}. L'aide auditive consomme un courant de \qty{1,2}{mA} et fonctionne 16 heures par jour.
Une batterie lithium-ion stocke environ \qty{250}{W.h} par kilogramme.
\end{docu}
\begin{donnees}
$F=\qty{96500}{C.mol^{-1}}$ ; $M(\ce{Zn})=\qty{65,4}{g.mol^{-1}}$ ; l'air contient 21~\% de dioxygène en volume ; volume molaire des gaz :
\qty{24}{L.mol^{-1}}.
\end{donnees}
\begin{questions}[resume]
\item Identifier l'anode et la cathode de la pile. Écrire l'équation de sa réaction de fonctionnement.
\item Exprimer la capacité de la pile en coulombs, puis calculer la masse minimale de zinc qu'elle contient.
\item Pendant combien de jours la pile peut-elle alimenter l'aide auditive ?
\item Calculer l'énergie électrique stockée, puis l'énergie stockée par kilogramme de pile. Comparer à une batterie
lithium-ion et conclure.
\item \initiative Quel volume d'air la pile doit-elle absorber au cours de son fonctionnement ? Pourquoi les trous de la
pile doivent-ils rester obturés jusqu'à la première utilisation ?
\end{questions}
\end{exoBac}

\begin{exoBac}{7}{Le saut en parachute}
Un parachutiste, de masse $m=\qty{80}{kg}$ avec son équipement, saute d'un avion. Pendant la chute libre, avant l'ouverture du
parachute, l'air exerce sur lui une force de frottement opposée à sa vitesse, de valeur $f=k\,v^2$, avec $k=\qty{0,25}{kg.m^{-1}}$.
On étudie son mouvement dans le référentiel terrestre supposé galiléen, sur un axe vertical orienté vers le bas ; la
poussée d'Archimède est négligée et $g=\qty{9,81}{m.s^{-2}}$.
\begin{questions}
\item Faire le bilan des forces qui s'exercent sur le parachutiste et les représenter sur un schéma.
\item Montrer que la vitesse vérifie l'équation différentielle $\dfrac{\mathrm{d}v}{\mathrm{d}t}=g-\dfrac{k}{m}\,v^2$.
\item Montrer que la vitesse tend vers une vitesse limite dont on donnera l'expression et la valeur, en \unit{m.s^{-1}} et en
\unit{km.h^{-1}}.
\item \capacite{Capacité numérique}Le programme suivant résout l'équation par la méthode d'Euler ; la figure représente la vitesse
obtenue.
\begin{codecote}[0.62]
\begin{pythoncol}
m, g, k = 80, 9.81, 0.25
vlim = (m * g / k) ** 0.5
t, v, z, dt = 0.0, 0.0, 0.0, 0.1
while v < 0.95 * vlim:
    v = v + (g - k / m * v**2) * dt
    z = z + v * dt
    t = t + dt
print(round(t, 1), round(v, 1), round(z))
\end{pythoncol}
\cote
\begin{consolecol}
10.4 53.3 373
\end{consolecol}
\end{codecote}
\schema{bac-s5-parachute}
Justifier la ligne \code{v = v + (g - k / m * v**2) * dt}. Que représentent les trois valeurs affichées ?
\item Vérifier, sur la figure, que l'accélération initiale vaut $g$. Justifier ce résultat.
\item Quelle distance le parachutiste aurait-il parcourue pendant la même durée s'il n'y avait pas de frottements ?
\item Lorsque la vitesse limite est atteinte, l'énergie cinétique reste constante. Que devient l'énergie potentielle de
pesanteur perdue ? Calculer la puissance correspondante.
\item \initiative Une fois le parachute ouvert, $k$ vaut \qty{30}{kg.m^{-1}}. Calculer la nouvelle vitesse limite. De quelle hauteur
faudrait-il sauter sans parachute pour toucher le sol à cette vitesse ?
\end{questions}
\end{exoBac}

\begin{exoBac}{5}{Le diamètre d'un cheveu}
Pour mesurer le diamètre $a$ d'un cheveu, on l'éclaire avec un laser : le cheveu diffracte la lumière comme le ferait une
fente de même largeur. On observe sur un écran une tache centrale de largeur $L$.
\schema{bac-s5-cheveu}
\begin{donnees}
$\lambda=\qty{650}{nm}$, d'incertitude négligeable ; $D=\qty{2,00}{m}$ avec $u(D)=\qty{0,01}{m}$ ; $L=\qty{3,7}{cm}$ avec $u(L)=\qty{0,1}{cm}$. Demi-angle
de la tache centrale : $\theta=\dfrac{\lambda}{a}$ ; pour les petits angles, $\tan\theta\approx\theta$. Incertitude :
$\dfrac{u(a)}{a}=\sqrt{\left(\dfrac{u(L)}{L}\right)^2+\left(\dfrac{u(D)}{D}\right)^2}$.
\end{donnees}
\begin{questions}
\item Pourquoi la diffraction par un cheveu est-elle facilement observable avec un laser ?
\item Montrer que $a=\dfrac{2\,\lambda\,D}{L}$.
\item Calculer le diamètre du cheveu.
\item Calculer l'incertitude-type sur $a$ et écrire le résultat de la mesure.
\item \initiative Un micromètre donne pour le même cheveu \qty{72}{\micro m}, avec une incertitude-type de \qty{1}{\micro m}. Les deux mesures
sont-elles compatibles ? On admet que l'incertitude-type sur l'écart des deux mesures vaut $\sqrt{u_1^2+u_2^2}$.
\item Proposer deux modifications du protocole qui réduiraient l'incertitude sur $a$.
\end{questions}
\end{exoBac}
